
\section{9. Discussion}

\subsection{9.1 Key Takeaways}

\subsubsection{9.1.1 Takeaway 1: LoRA is a Strong Distillation Target}

Training only 1\% of parameters (LoRA) recovers most of the accuracy lost 
by head-dim reduction.

\subsubsection{9.1.2 Takeaway 2: Head-Dim Reduction Works}

Reducing head dimension by 2$\times$ gives ~{{PARAM\_REDUCTION}}$\times$ parameter 
reduction with acceptable loss.

\subsubsection{9.1.3 Takeaway 3: Drop-In Wrapper}

Squeezing Heads preserves pretrained weights exactly. It can be applied 
to any model without retraining.

\subsection{9.2 When to Use Squeezing Heads}

\subsubsection{Recommended}

\begin{itemize}
\item \textbf{Model compression} where head pruning loses too much.
\end{itemize}
\begin{itemize}
\item \textbf{Continual learning} where you can't retrain from scratch.
\end{itemize}
\begin{itemize}
\item \textbf{Multi-task adaptation} -- one LoRA per task.
\end{itemize}
\begin{itemize}
\item \textbf{Research on head structure} -- study what's lost by dim reduction.
\end{itemize}

\subsubsection{Not Recommended}

\begin{itemize}
\item \textbf{Inference speedup} -- SH adds LoRA overhead.
\end{itemize}
\begin{itemize}
\item \textbf{Very small models} -- LoRA overhead not worth it.
\end{itemize}
\begin{itemize}
\item \textbf{Training from scratch} -- LoRA is a fine-tuning technique.
\end{itemize}

\subsection{9.3 Choosing LoRA Rank}

\begin{table}[H]
\centering
\begin{tabular}{ll}
\toprule
\textbf{Use case} & \textbf{Recommended $r_L$} \\
\midrule
Extreme compression & 4 \\
Balanced & 8 \\
Safe & 16 \\
High quality & 32 \\
\bottomrule
\end{tabular}
\end{table}

\subsection{9.4 Choosing Head-Dim Reduction}

\begin{table}[H]
\centering
\begin{tabular}{ll}
\toprule
\textbf{Use case} & \textbf{Recommended $s$} \\
\midrule
Safe & 1.0 (no reduction) \\
Balanced & 1.5 \\
Aggressive & 2.0 \\
Extreme & 2.5 \\
\bottomrule
\end{tabular}
\end{table}

\subsection{9.5 Cost-Benefit}

For DistilBERT at $L = 128$:

\begin{table}[H]
\centering
\begin{tabular}{llll}
\toprule
\textbf{Config} & \textbf{Trainable} & \textbf{Acc} & \textbf{$\Delta$} \\
\midrule
Full FT & 66.9M & 90.71\% & -- \\
SH $r_L=16$ & {{D\_LORA}} & {{D\_ACC}}\% & {{D\_DELTA}}\% \\
SH + $s=2$ & {{D\_LORA}} & {{D\_S20}}\% & {{D\_S20\_D}}\% \\
\bottomrule
\end{tabular}
\end{table}

\subsection{9.6 Open Questions}

1. \textbf{Optimal per-head rank:} Can we allocate LoRA rank per head?
2. \textbf{Combining with pruning:} SH + head pruning?
3. \textbf{Distillation from teacher:} Can SH learn from a larger teacher?
4. \textbf{Cross-task transfer:} Does one LoRA transfer to another task?

\subsection{9.7 Broader Impact}

\subsubsection{Positive}
\begin{itemize}
\item \textbf{Cheaper fine-tuning:} LoRA trains 1\% of parameters.
\end{itemize}
\begin{itemize}
\item \textbf{Preserves pretrained:} No catastrophic forgetting.
\end{itemize}
\begin{itemize}
\item \textbf{Multi-task:} One adapter per task.
\end{itemize}

\subsubsection{Negative}
\begin{itemize}
\item \textbf{Concentration:} Large labs benefit most.
\end{itemize}
\begin{itemize}
\item \textbf{Rebound effect:} Cheaper fine-tuning may increase total compute.
\end{itemize}

\subsubsection{Environmental}
Total compute: ~{{COMPUTE\_HOURS}} GPU-hours, ~{{CO2}} kg CO$_2$eq.